
kriyajnana

I broke the Sankhya into two main parts.
Writing this I saw that there is a division indicated here
that might be important: the division between
apara and para indriya jaya.

ekagrata

sarva arthata ekagrata ksayodayau cittasya samadhi-parinama

$19

The logical form of all judgments consists of the objective unity of the apperception of the concepts contained therein

buddhi -- ahamkara -- manas

All our cognitions
Begin with Grahya
Proceed through Grahana
End with Grahitri

samapatti
[A U M]

Sankhya says the inner instrument is known through its functions:


The Concept as such

Manas, the understanding, first the faculty for the cognition of the general (of rules)
this is Nirodha parinama

The understanding works with pratyahara which returns a manifold of sensation and
the understanding converts this to a representation.
I think of this as a "Closure/Dharma chakra".
It is consuming perceptions into a concept of an object.
The inner instrument as a machine is a kantian thing I admit.
But extracting the bija from this inner instrument is the goal.

....Bill....


Judgment

Ahamkara, the determinative power of judgment, second the faculty for the subsumption of the particular under the general
this is Samadhi parinama

....Bill is a man...

Syllogism

Buddhi, reason, third the faculty for the determination of the particular from the general (for the derivation from principles)
this is Ekagrata parinama

A Change of Ekagrata
from Bill is a man to Bill is mortal
based on the same representation ....Bill qua man....

This is a Determinate Power of judgment
as opposed to a Reflective Power of judgment.
It is an Empirical Science in action ...this is Kriya Jnana
Rajo-guna Middle Way through and through
Samadhi middle way between Prajna and Dharma padas.

Yoga sutra is a grammar. It is the seed form of sarvajna.
The seed form of "the seed". So it is what it is.
It is definitely presenting the three functions of the inner instrument.
In its full specification it is Tarka.

This is Patanjali's definition of Jnana as part of the definition of Buddhi tattva

samyama is jnana from prajna (philosophical cognition is "derivation from principles")


prajna
kriya
jnana
dharma

prajna : viraga or aishvarya : samapatti
jnana  : samadhi or pratibha : pratipatti
